\chapter{Literature Review}
\label{chap:lr}
\chaptermark{Literature Review}

This chapter presents a focused narrative review of research on bug report quality, interactive reporting, and language-model-assisted report improvement. It also examines the use of external documentation and the Model Context Protocol (MCP), which are relevant to the proposed system. The review distinguishes evidence about report content from evidence about the reporting interaction and subsequent debugging. It does not claim exhaustive coverage or constitute a systematic literature review.

\section{The Role of Bug Reports in Software Maintenance}

Bug reports connect a user's account of a software problem with the information needed by developers to investigate it. Zimmermann et al. surveyed developers and users of Apache, Eclipse, and Mozilla and identified a mismatch between the information developers need and what reporters can provide. Reproduction steps, stack traces, and test cases were considered useful but were also difficult for reporters to supply \cite{zimmermann2010}. This finding motivates reporting assistance that helps people communicate evidence they possess, while recognizing that they may not have access to diagnostic artifacts.

Soltani et al. examined the importance of report elements through interviews with 35 developers, a survey of 305 developers, and an analysis of issue repositories from 250 projects. Their results emphasize reproduction steps and stack traces, while also considering software versions, user content, and other contextual information \cite{soltani2020}. The study provides broader evidence that useful report content extends beyond a general symptom description. Its repository analysis concerns associations between report elements and resolution times; such associations do not establish that adding a field to a reporting interface will cause faster resolution.

Wang et al. compared initial and final versions of 1,000 bug reports from ten open-source web applications. They found that end-user usage information and developers' diagnostic suggestions were often absent initially and added during later discussions \cite{wang2025}. This report evolution suggests an opportunity to collect relevant usage context earlier. However, diagnostic suggestions added during investigation should not be confused with facts that the original reporter could have supplied.

Erfani Joorabchi et al. studied non-reproducible reports in one industrial and five open-source repositories. Their analysis identified several reasons for non-reproducibility, including environmental differences, insufficient information, and dependencies between bugs \cite{joorabchi2014}. Consequently, a clear report can support investigation without ensuring successful reproduction: some obstacles concern the software and its environment rather than the quality of the description.

Together, these studies justify improving information collection at reporting time. They also limit the conclusions that can be drawn from report quality alone. A reporting system can improve the evidence available for investigation, but its effect on defect localization or repair must be demonstrated separately.

\section{Completeness, Clarity, and Factual Accuracy}

Chaparro et al. analyzed 2,912 reports and investigated the detection of missing expected behavior and steps to reproduce. Their findings show that these elements are frequently absent, even when observed behavior is described \cite{chaparro2017}. This distinction is important because describing a failure does not necessarily explain either the intended outcome or the actions that preceded it.

The following information categories synthesize the report elements investigated in the literature \cite{zimmermann2010,chaparro2017,soltani2020,wang2025}:

\begin{itemize}
    \item \textbf{Observed behavior:} what happened, including the visible symptom or available error message.
    \item \textbf{Expected behavior:} the outcome the reporter expected, with its basis clarified where necessary.
    \item \textbf{Steps to reproduce:} the actions, inputs, and relevant preconditions preceding the problem.
    \item \textbf{Context:} applicable details about the service, version, environment, and circumstances of use.
    \item \textbf{Supporting evidence:} screenshots, logs, stack traces, or other artifacts when available and relevant.
\end{itemize}

Completeness is only one aspect of quality. Chaparro et al.'s EULER approach assesses reproduction steps by relating natural-language descriptions to application interactions and producing feedback about problematic or missing steps \cite{chaparro2019}. A report can therefore contain a reproduction section while still omitting necessary actions or describing them ambiguously. Counting completed fields alone cannot establish whether another person can follow the report.

Factual accuracy creates a further distinction for generated reports. Ji et al. describe hallucination in natural-language generation as output that is unsupported by, or inconsistent with, the relevant source information \cite{ji2023}. Applied to reporting, a fluent description may still contain an invented error message, an unconfirmed action, or an altered observation. These examples are consequences of the distinction between linguistic plausibility and evidential support, rather than findings about the proposed prototype.

For this thesis, completeness refers to the presence of relevant information, clarity to its comprehensibility, and factual accuracy to its consistency with available evidence. A field that is inapplicable or genuinely unknown should remain distinguishable from an omission. These are analytical definitions used to organize the review; their operational assessment belongs to the evaluation methodology.

\section{Structured Forms and Interactive Reporting}

Structured templates organize information into predefined categories. The template-based reporting interface used as a comparison in the BURT study illustrates this approach: reporters supply information through designated sections rather than receiving an adaptive dialogue \cite{song2022}. Structure supplies an organizational framework, but the effectiveness of additional guidance remains an empirical question.

CUEZILLA, described by Zimmermann et al., assesses report quality and recommends additional report elements. It was trained and evaluated using 289 reports rated by developers \cite{zimmermann2010}. It establishes that feedback during reporting predates LLM-based systems. In this review, the journal article provides the principal account of that work; the earlier conference publication is not treated as an independent replication.

Song et al. introduced BURT, a task-oriented chatbot that guides end-users through observed behavior, expected behavior, and reproduction steps. It provides immediate quality feedback, clarification prompts, and graphical suggestions. Its implemented reporting interface is web-based, while its application analysis and evaluation concern Android apps \cite{song2022}. These are different platform properties: a browser-based interface does not make the evaluated reporting technique a method for arbitrary web applications.

BURT combines language processing with a graph of application screens and interactions derived from execution data. Its evaluation used twelve bugs from six Android apps, compared reports with those from a template-based system, and examined users' experiences and the correctness of reported information \cite{song2022}. The results support interactive guidance in the studied setting. They do not establish equivalent effects for every service or documentation source. Importantly, interaction and usability are already part of this line of work, so considering user effort is not by itself a new contribution.

The relevant distinction is therefore between different mechanisms for guidance and different sources of application context. The proposed system will use service documentation accessed through MCP to inform clarification questions. Whether this approach produces useful questions and preserves the reporter's account requires evaluation; the existence of an interactive interface alone does not distinguish it from BURT.

\section{Natural Language Processing and Large Language Models}

The missing-information detection approach of Chaparro et al. demonstrates that bug descriptions can be analyzed automatically for particular report elements \cite{chaparro2017}. Detection, however, answers a narrower question than conversational assistance. Recognizing an absent reproduction description does not directly establish which question should be asked next or which service-specific precondition matters.

Hou et al.'s systematic literature review describes applications of LLMs across software engineering and discusses challenges in their use and evaluation \cite{hou2024}. For reporting, language models offer a way to interpret varied descriptions and produce questions or summaries. This is a rationale for investigating them, rather than evidence that an LLM necessarily outperforms a specialized classifier or a carefully designed form.

McGowan et al. describe work on using application interface and execution information to prompt an LLM to improve reports and predict reproduction steps \cite{mcgowan2023}. Their publication states that evaluation of the tool was ongoing. It therefore provides an example of an application-informed design, but should not be cited as a completed demonstration of improved reporting performance.

Hallucination is a material limitation when generation is used to complete an account of an actual event \cite{ji2023}. In the proposed system, a question suggested by documentation will be treated as a request for confirmation. A documented upload limit, for example, can motivate asking about file size; it cannot establish the size of the file the user actually submitted. This example illustrates the intended separation between documented behavior, user observations, and unresolved information.

The role investigated here is assistance with problem description and information collection. Identifying what the user observed and under which conditions is different from determining the underlying code defect. The literature on report analysis provides a foundation for the former, without removing the need for independent evidence about the latter.

\section{LLM-Based Bug Report Improvement}

Bo et al. introduced ChatBR, which combines report-element assessment with ChatGPT-based improvement. The approach addresses observed behavior, expected behavior, and reproduction steps, and evaluates missing-information detection and generated content \cite{chatbr2024}. This is relevant to identifying deficiencies and producing structured reports. However, similarity to a reference report is a different outcome from obtaining confirmation from the person who experienced a particular failure.

Acharya and Ginde evaluate instruction-tuned models and ChatGPT for converting informal reports into structured reports. Their study uses quality and text-similarity measures and compares performance across models and report elements \cite{acharya2025}. The published EASE 2025 paper is used here rather than its earlier arXiv version. Its results concern report generation under the studied conditions; they do not directly measure an interactive reporter's workload.

Mahmud et al. introduced AstroBR, which uses an LLM together with a dynamically constructed application model to identify, extract, and assess reproduction steps \cite{mahmud2025}. Its evaluation compares step annotations and missing-step suggestions with reference annotations and EULER. The approach demonstrates that LLM-based reporting assistance can use application-specific evidence beyond the report text. Its reliance on GUI execution models distinguishes its context source from the service documentation proposed here, without implying that either source is universally superior.

Mukherjee and Rahman's BugMentor retrieves relevant historical reports and uses neural text generation to answer follow-up questions about an existing report. The evaluation includes reference-answer metrics and a developer study \cite{bugmentor2026}. BugMentor thus addresses an important part of the clarification process, but generating an answer from historical context differs from asking the current reporter to supply an observation. This distinction concerns the origin of the information, not simply whether retrieval is used.

Akyol et al.'s ImproBR combines detection of incomplete or ambiguous report sections with LLM-based improvement and retrieval from Minecraft Wiki. It is listed in the EASE 2026 research programme; the accessible author manuscript is cited explicitly \cite{akyol2026}. Its evaluation includes structural completeness, executability, and reproducibility of Minecraft reports. This is a particularly close example of report improvement using external domain knowledge. Accordingly, combining bug reporting with retrieved documentation cannot itself be claimed as an unexplored idea. The present thesis instead investigates documentation-informed questions within an interaction with the reporter.

Chen et al. investigate decomposing complex bug reports using LLMs and show that performance depends on prompting \cite{chen2025}. Their arXiv manuscript is cited as a preprint. Although decomposition differs from information collection, it provides a related example of the need to evaluate the interpretation of report content rather than assuming correct understanding.

Across these approaches, three operations should be distinguished: detecting a deficiency, generating a candidate completion, and collecting confirmed information. The proposed system focuses on the third operation while using the first to guide the interaction. It may structure the answers, but any proposed completion remains subject to confirmation.

\section{Service Documentation and MCP}

Lewis et al. introduced retrieval-augmented generation (RAG), combining a generative model with information retrieved from an external collection. Their experiments concern knowledge-intensive NLP tasks and show benefits of retrieval in those settings \cite{lewis2020}. The general relevance to this thesis is the use of identifiable external information rather than reliance on model parameters alone. Those experiments do not directly establish the quality of documentation-informed bug-reporting questions.

External context is already represented in the reviewed reporting literature: AstroBR uses an application execution model, BugMentor retrieves historical reports, and ImproBR retrieves domain knowledge \cite{mahmud2025,bugmentor2026,akyol2026}. These sources answer different questions. Documentation can describe supported operations and prerequisites, while an execution trace or the reporter's answer can supply evidence about what actually happened.

MCP standardizes communication between applications and servers exposing context and capabilities. Its specification defines resources for application-managed contextual data and tools for callable operations \cite{mcp2025resources,mcp2025tools}. A documentation integration may expose documents as resources or provide search and retrieval tools. MCP does not itself choose the relevant passages, establish their correctness, or guarantee that a model's output is grounded. It is the access mechanism; retrieval and question-generation policies remain separate system responsibilities.

The intended role of MCP in this thesis is to provide access to documentation of the services about which users report problems. The system will use that material to formulate questions about relevant operations, inputs, and preconditions, with the purpose of making the reported problem more precisely specified. This is a design commitment, not a result established by the reviewed literature. It does not require assuming a particular service, implementation stack, or industrial partner before those details are defined.

Access to more text also does not ensure effective use of that text. Liu et al. found that language-model performance on long-context tasks can depend on the position of relevant information, with poorer performance when it appears in the middle \cite{liu2024}. This finding motivates selecting relevant documentation rather than treating a large context window as sufficient. Together with the hallucination literature, it supports assessing whether the questions and summaries are actually supported by the retrieved material \cite{ji2023}.

For the proposed reporting interaction, documentation will inform what to ask, while the user's answers will establish their reported experience. A documented expected outcome may be presented as documentation-derived information; it must not silently replace the user's own expectation or observation. Where the two disagree, the disagreement is information to preserve and clarify.

\section{Evaluation Practices in Prior Research}

The reviewed studies evaluate different objects, and their results are not interchangeable. Chaparro et al. assess detection of missing information, while EULER and AstroBR evaluate reproduction-step quality against annotations \cite{chaparro2017,chaparro2019,mahmud2025}. Such measures assess an analysis component rather than the whole reporting experience. They do not on their own measure whether a user can complete an interaction successfully.

ChatBR and Acharya and Ginde use measures of generated report quality and similarity to reference content \cite{chatbr2024,acharya2025}. Reference-based comparison supports reproducible assessment, but a generated completion can resemble a reference without being confirmed for the current incident. ImproBR extends evaluation to executability and reproducibility \cite{akyol2026}; this illustrates the value of distinguishing a well-formed report from one whose steps can actually be followed.

BURT combines report-content evaluation with an investigation of user experience \cite{song2022}. This is directly relevant to an interactive form because users both supply the information and respond to the system's guidance. The additional questions may improve a report while creating additional interaction. Quality and reporting effort should therefore be treated as separate outcomes rather than combined into an unspecified claim of efficiency.

Subjective workload is also distinct from elapsed time. Hart and Staveland's NASA-TLX assesses dimensions including mental demand, temporal demand, effort, frustration, and perceived performance \cite{hart1988}. It provides an established basis for considering perceived workload alongside behavioral measurements. Its inclusion in the review does not establish which instrument or scoring procedure will be adopted; those choices require specification in the evaluation chapter.

Finally, report preparation, understanding a report, reproducing a failure, and repairing a defect are different stages. Repository studies describe relationships between report content and maintenance outcomes \cite{soltani2020,joorabchi2014}, but those relationships do not license a causal claim about the proposed interface. An evaluation of report quality and reporting effort can support conclusions about bug reporting. Claims about faster diagnosis or resolution would require measurements of those downstream activities.

\section{Synthesis, Research Gap, and Research Questions}

The literature establishes that incomplete and ambiguous information is a recurring reporting problem. It also provides several existing responses: quality feedback, interactive dialogue, application-informed assessment, LLM-based rewriting, and retrieval-assisted generation. Table~\ref{tab:reporting-approaches} summarizes selected approaches by their main task and source of context. The categories describe the cited systems rather than an exhaustive taxonomy.

\begingroup
\small
\setlength{\tabcolsep}{4pt}
\renewcommand{\arraystretch}{1.15}
\begin{longtable}{@{}>{\raggedright\arraybackslash}p{0.22\textwidth}>{\raggedright\arraybackslash}p{0.34\textwidth}>{\raggedright\arraybackslash}p{0.38\textwidth}@{}}
\caption{Selected approaches relevant to intelligent bug reporting.}\label{tab:reporting-approaches}\\
\toprule
\textbf{Approach} & \textbf{Main task} & \textbf{Context or evidence} \\
\midrule
\endfirsthead
\toprule
\textbf{Approach} & \textbf{Main task} & \textbf{Context or evidence} \\
\midrule
\endhead
\bottomrule
\endfoot
CUEZILLA \cite{zimmermann2010} & Assess reports and recommend additional elements & Report features and developer quality ratings \\
EULER \cite{chaparro2019} & Assess reproduction steps and provide feedback & Report text and application interactions \\
BURT \cite{song2022} & Guide reporters through an interactive dialogue & Reporter answers and an Android application model \\
ChatBR \cite{chatbr2024} & Assess and improve report elements & Report text and generated completions \\
AstroBR \cite{mahmud2025} & Assess steps and suggest missing steps & Report text and a GUI execution model \\
BugMentor \cite{bugmentor2026} & Generate answers to follow-up questions & Current and retrieved historical reports \\
ImproBR \cite{akyol2026} & Improve incomplete or ambiguous sections & Report text and retrieved domain knowledge \\
Proposed system & Ask clarification questions and structure confirmed answers & Reporter input and service documentation accessed through MCP \\
\end{longtable}
\endgroup

The research opportunity is not the invention of interactive reporting or the first use of external knowledge. The reviewed work motivates a more specific empirical question: how useful is a documentation-informed LLM reporting interaction when assessed jointly for report quality, reporting effort, and fidelity to the reporter's evidence? The use of MCP identifies how documentation will be accessed; its presence alone is not evidence of improved reporting.

The intended organizational clients are companies integrating the prototype into support processes for their software services. The reporters are users of those services, and developers and support specialists receive the resulting reports. The project is independently developed and has no industrial partner; its evaluation will use controlled support scenarios. The planned documentation collection describes the adopting company's service capabilities and previously resolved problems. The intended interaction allows users to write in their own language and receive clarification questions in that language, and produces a structured report for developers. These are project requirements rather than evidence of equivalent performance across languages. The languages included in the evaluation and the provenance of example documentation and resolved cases must be specified in the methodology.

The intended contribution is a prototype and an evaluation of this combination in that service context. The evaluation should separate the effect of interactive assistance from the additional contribution of documentation access. A conventional form provides a reference for reporting practice, while an otherwise comparable assisted form without documentation would help isolate the value of that context. This distinction identifies an attribution issue arising from the review; the experimental protocol is to be specified in Chapter~5.

This synthesis motivates the four research questions stated in Section~\ref{sec:research-questions}. RQ1 addresses completeness and clarity, reflecting the distinction between the presence of report elements and their usefulness. RQ2 addresses reporting effort and workload, recognizing that additional guidance may also demand more from the reporter. RQ3 examines fidelity to user-confirmed information, which cannot be established by fluency or structural completeness alone. RQ4 examines the additional effect of documentation, distinguishing it from the effect of LLM-assisted interaction. Together, these questions concern the preparation of a useful and faithful problem description; they do not assume that the system can identify the root cause or accelerate software repair.
